% Part: methods
% Chapter: proofs
% Section: starting-proofs

\documentclass[../../../include/open-logic-section]{subfiles}

\begin{document}

\olfileid{mth}{prf}{str}

\olsection{如何开始证明}

但是，究竟该从哪里开始呢？

你拿到了需要证明的命题，因此它应当是证明中最后才出现的内容（当然，你可以在开头\emph{宣告}你要证明它，但不能把它当作假设来使用）。在一张空白纸的最下方写下你要证明的结论——这样你就不会偏离目标。

接下来，你可能有一些可以使用的假设（当我们在下一节讨论你所进行的证明的\emph{类型}时，这一点会更加清楚）。把这些假设写在页面的顶部，并务必明确标注它们是假设（即，如果你假设 $p$，就写上“假设 $p$”或“设 $p$”）。最后，题目中可能有一些你需要知晓的定义。你可能被要求使用某个特定的定义，或者在你所要处理的假设或结论中可能包含各种定义。\emph{把它们写下来，并确保你理解它们的含义。}

你构建证明的方式也取决于题目的形式。下一节将根据语句的类型，详细说明如何构建你的证明。

\end{document}